\documentclass[11pt]{article}

% --------------------------------------------------
% Packages
% --------------------------------------------------
\usepackage[T1]{fontenc}
\usepackage{lmodern}
\usepackage{geometry}
\usepackage{microtype}
\usepackage{setspace}
\usepackage{amsmath,amssymb,amsthm,mathtools}
\usepackage{csquotes}
\usepackage{hyperref}
\usepackage{booktabs}

\geometry{margin=1in}
\setstretch{1.15}

% --------------------------------------------------
% Theorem Environments
% --------------------------------------------------
\newtheorem{definition}{Definition}
\newtheorem{assumption}{Assumption}
\newtheorem{lemma}{Lemma}
\newtheorem{proposition}{Proposition}
\newtheorem{theorem}{Theorem}
\theoremstyle{remark}
\newtheorem{remark}{Remark}

% --------------------------------------------------
% Macros
% --------------------------------------------------
\newcommand{\doop}{\mathrm{do}}
\newcommand{\Adm}{\mathrm{Adm}}
\newcommand{\Hist}{\mathcal{H}}
\newcommand{\Fut}{\mathcal{F}}
\newcommand{\Val}{\mathsf{Val}}
\newcommand{\Reg}{\mathsf{Reg}}
\newcommand{\Tnode}{\mathbf{T}}
\newcommand{\Enode}{\mathbf{E}}
\newcommand{\State}{\mathcal{S}}
\newcommand{\Event}{\mathcal{E}}
\newcommand{\Future}{\mathcal{F}}
\newcommand{\Noise}{\mathcal{N}}
\newcommand{\Graph}{\mathcal{G}}
\newcommand{\NN}{\mathcal{N}}
\newcommand{\Alg}{\mathcal{A}}
\newcommand{\Do}{\mathrm{do}}
\newcommand{\Var}{\mathcal{V}}

% --------------------------------------------------
% Title
% --------------------------------------------------
\title{Causal Mechanistic Interpretability: From Inspection to Event-Historical Structures, Two-Channel Mediation, and Resilience in Neural and Knowledge Systems}
\author{Flyxion}
\date{December 2025}

\begin{document}
\maketitle

% ==================================================
\begin{abstract}
% ==================================================

Recent advances in mechanistic interpretability have grounded explanations of neural networks in causal reasoning, progressing from descriptive inspection to interventions, counterfactuals, and causal abstractions. However, state-centric approaches overlook irreversibility, history sensitivity, and regime shifts. This essay synthesizes a unified causal paradigm, introducing a typed framework distinguishing \emph{texture-nodes} (value-carrying variables) from \emph{event-nodes} (regime-selecting constraints). We formalize two-channel mediation, graph rewrite semantics, and event-triggered irreversibility, proving decomposition theorems that clarify substitution failures and algorithmic mixtures. The framework extends to knowledge networks, where resilience depends on event-level structure. By embedding history within causal graphs, mechanistic understanding becomes temporal, plural, and constraint-aware, resolving tensions in existing methods while preserving empirical rigor.

\end{abstract}

\tableofcontents

% ==================================================
\section{Introduction}
% ==================================================

The remarkable capabilities of large neural networks have spurred efforts to understand their internal operations. Early interpretability work emphasized visualization, feature attribution, and activation patterns, but these observational methods often conflate correlation with computation. Mechanistic interpretability seeks to move beyond such limitations by identifying the causal processes that generate behavior.

A causal paradigm has emerged, defining understanding as the ability to intervene on internal variables, evaluate counterfactuals, and validate algorithmic abstractions. This view treats neural networks as causal systems amenable to mediation analysis, tracing, and substitution experiments. Yet disagreement persists about what constitutes genuine mechanism, with definitions oscillating between cultural motives, technical access, and epistemic standards.

This essay synthesizes and extends this paradigm. We trace its progression from activation steering to causal abstraction, highlight state-centric limitations, and introduce a minimal typed extension with \emph{event-nodes} and \emph{texture-nodes}. This distinction enables two-channel mediation, graph rewrites, and history-sensitive explanations, unifying disparate techniques while clarifying irreversibility and regime shifts. We further extend the framework to knowledge networks, analyzing resilience under perturbation.

The result is a rigorous, temporal account of mechanistic understanding as causal reverse engineering grounded in irreversible events.

% ==================================================
\section{What Does ``Mechanistic'' Mean?}
% ==================================================

The term \emph{mechanistic} is ambiguous. Cultural definitions emphasize community or motivation, such as AI safety concerns or curiosity about black boxes. Technical definitions vary from broad inspection of internals to narrow causal explanation.

We adopt a demanding technical sense: a system is mechanistically understood if internal components can be identified whose interventions produce predictable changes consistent with an algorithmic description. This excludes pure observation; correlation does not imply control.

Causality distinguishes mechanisms from regularities. Interventions set variables independently of antecedents, enabling tests of necessity and sufficiency. Counterfactuals extend this to hypothetical configurations.

Without causality, interpretability risks pattern cataloging. With it, it becomes a science of explanation and control.

% ==================================================
\section{Causality as the Foundation of Mechanism}
% ==================================================

\begin{definition}[Intervention]
An intervention on a neural network sets an internal variable to a value it would not naturally take, breaking upstream dependencies.
\end{definition}

Causality evaluates how systems respond to such changes. A component is relevant if altering it systematically affects behavior.

This foundation supports progression from linear steering to mediation and abstraction.

% ==================================================
\section{Activation Steering as Linear Intervention}
% ==================================================

Activation steering adds vectors to internals, biasing behaviors. Steering vectors, derived from activation differences, reveal linear directions in representation space.

While demonstrating causal connections, steering is brittle and does not identify algorithms. Its value lies in establishing representations as non-epiphenomenal.

% ==================================================
\section{Causal Mediation and Tracing Information Flow}
% ==================================================

Causal mediation decomposes input-output effects through intermediates. In networks, tracing corrupts inputs and restores components, identifying mediators whose recovery restores behavior.

This reveals pipelines: feedforward layers retrieve, attention routes, later stages assemble. Claims rest on intervention, grounding information flow causally.

% ==================================================
\section{Causal Abstraction and Algorithmic Understanding}
% ==================================================

Causal abstraction relates networks to algorithms via interventional equivalence.

\begin{definition}[Causal Abstraction]
A high-level model $\Alg$ abstracts a low-level system if mappings exist such that interventions produce equivalent outputs.
\end{definition}

Interchange interventions substitute values from alternate inputs, validating abstractions if outputs match algorithmic predictions.

This is stringent, requiring not similarity but equivalence under substitution.

% ==================================================
\section{Limits of State-Based Approaches}
% ==================================================

State-based models privilege present values, treating history as dispensable. However, identical states reached differently may diverge under intervention due to irreversible commitments.

Functionalism collapses such distinctions, ignoring entropy and history. Causal explanation requires embedding function in event-historical frameworks.

Interpretability thus becomes historical reconstruction: inferring irreversible sequences from structure.

% ==================================================
\section{Event-Nodes and Texture-Nodes}
% ==================================================

To address these limits, we introduce typed causal nodes.

\begin{definition}[Texture-Node]
A texture-node carries a value $X_v \in \Val_v$; interventions substitute values.
\end{definition}

\begin{definition}[Event-Node]
An event-node selects a regime $E_v \in \Reg_v$, constraining admissible futures $\Adm(\Hist_t)$.
\end{definition}

Texture-nodes modulate within regimes; event-nodes select generators, inducing irreversibility.

This parallels procedural generation: parameters (textures) vs. generator choice (events).

% ==================================================
\section{Typed Causal Graphs and Rewrite Semantics}
% ==================================================

A typed model is $(G, \tau, \Val, \Reg, \Adm)$, with $\tau: V \to \{\Tnode, \Enode\}$.

Event interventions trigger rewrites, replacing subgraphs while preserving interfaces. Irreversibility arises structurally, not from loss.

% ==================================================
\section{Two-Channel Causal Mediation}
% ==================================================

Mediation decomposes along channels.

\begin{theorem}[Two-Channel Mediation]
The total effect decomposes into texture- and event-mediated components. Texture suffices iff events are invariant.
\end{theorem}

Proof sketches via nested counterfactuals in typed models.

This clarifies tracing: restoration probes both channels, explaining block structures and regime shifts.

% ==================================================
\section{Irreversibility, History, and Responsibility}
% ==================================================

Event-nodes formalize structural irreversibility: interventions restrict futures undonably by textures.

History matters as constraint; responsibility includes suppression of alternatives.

% ==================================================
\section{Extension to Knowledge Networks and Resilience}
% ==================================================

In knowledge systems (e.g., wikis), articles/links are textures; policies/bans are events.

Resilience requires texture redundancy and event stability. Perturbing events alters evolvability irreversibly.

% ==================================================
\section{Implications and Formal Results}
% ==================================================

Key propositions:
- Texture interventions reversible; events irreversible.
- Interchange valid iff event normal forms align.
- Abstractions must preserve irreversibility.

Understanding is reconstruction of values and constraints under intervention.

% ==================================================
\section{Against State-Based Functionalism}
% ==================================================

Functionalism erases history, collapsing causal equivalence. RSVP ontology unifies causality, thermodynamics, representation.

% ==================================================
\section{Conclusion}
% ==================================================

This synthesis reframes mechanistic interpretability as event-historical causal reverse engineering. By distinguishing nodes and channels, we resolve ambiguities, unify methods, and extend to broader systems. Understanding emerges from interventions respecting irreversible constraints.

\end{document}